\documentclass{Ilaris}
%\usepackage{glossary-mcols}
\usepackage[section=chapter]{glossaries}

\renewcommand*{\glsclearpage}{}
% glossar laden
\makeglossaries
\loadglsentries{glossar}

\newcommand{\gl}[1]{
    \underline{\gls{#1}}
}

\newcommand{\glossar}{
    \printglossary[title=Glossar]
}
\newcommand{\komplettesglossar}{
    \glsaddall
    \printglossaries
}

\newcommand{\glossareintrag}[2]{\newglossaryentry{#1}{#2}}

% zweispaltig
\renewcommand*{\glossarysection}[2][]{%
    \twocolumn[{\kapitel{#2}}]%
    %\twocolumn%
    \setlength\glsdescwidth{0.6\linewidth}%
    \glsglossarymark{\glossarytoctitle}%
}
\renewcommand*{\glossarypostamble}{\onecolumn}


\newglossaryentry{geschichte}
{
    name={\textbackslash geschichte},
    description={
            Fluff, Liedtexte, Zitate und vieles mehr können in dem kursiven, eingerückten Geschichtsblock stehen.\newline
            \textbf{Beispiel: }\textbackslash geschichte\{Der Text...\}\{Die Autorin\}
        }
}
\newglossaryentry{anf}
{
    name={\textbackslash anf},
    description={
            Setzt einen Text in Wörtlicherede. Verwendet deutsche Anführungszeichen unten und oben.
            \textbf{Beispiel: }\textbackslash anf\{Hey! Ich rede mit Euch.\}
        }
}
\newglossaryentry{kapitel}
{
    name={\textbackslash kapitel},
    description={
            Leitet Kapitel mit der angegeben Überschrift ein und registriert sie automatisch im Inhaltsverzeichnis. Innerhalb eines Zweispaltenlayouts am besten \textbackslash skapitel verwenden.
            \textbf{Beispiel}: \textbackslash kapitel\{Anhang\}}
}

\newglossaryentry{spaltenanfang}
{
    name={\textbackslash spaltenanfang},
    description={
            Beginnt das zweispaltige Layout bis es mit \textbackslash spaltenende abgeschlossen wird. Die Befehle müssen paarweise auftreten, sonst kommt Latex durcheinander.\\
            \textbf{Beispiel}: \textbackslash spaltenanfang}
}
\newglossaryentry{skapitel}
{
    name={\textbackslash kapitel},
    description={
            Kurzform für spaltenübergreifende Kapitel. Nachdem das Spaltenlayout mit \textbackslash spaltenanfang begonnen wurde, sollte dieser Befehl verwendet werden, der automatisch dafür sorgt, dass die Kapitelüberschrift über beide Spalten geht und der folgende Text wieder im Spaltenlayout steht.\\
            \textbf{Beispiel}: \textbackslash skapitel\{Einleitung\}
        }
}
\newglossaryentry{abschnitt}
{
    name={\textbackslash abschnitt},
    description={
            Abschnitte (Unterkapitel) können ein Kapitel unterteilen und stehen ebenfalls im Inhaltsverzeichnis. Die Abschnittsüberschriften sollen das Spaltenlayout (falls verwendet) nicht unterbrechen.\newline
            \textbf{Beispiel}: \textbackslash abschnitt\{Die Ankunft\}
        }
}
\newglossaryentry{absatz}
{
    name={\textbackslash absatz},
    description={
            Die dritte Stufe der Inhaltlichen Strukturierung. Absätze können mit diesem Befehl noch eine kleine Überschrift bekommen, die allerdings nichtmehr im Inhaltsverzeichnis auftaucht. Sie hebt sich visuell aber noch deutlich ab um so Absätze im Text auf den ersten Blick zu finden. In Abenteuern eigenen sie sich zum Beispiel um auf Beschreibungen bestimmter Orte, Scenen oder Personen hinzuweisen.
            \newline
            \textbf{Beispiel}: \textbackslash absatz\{Der Geheimgang\}
        }
}

\newglossaryentry{neuezeile}
{
    name={\textbackslash neuezeile},
    description={
            Dies ist nichts weiter als ein Zeilenumbruch, der nur der vollständigkeit halber einen eigenen Befehl bekommen hat. Er leitet direkt auf \textbackslash newline weiter. \newline
            \textbf{Beispiel}: \textbackslash neuezeile
        }
}
\newglossaryentry{neuespalte}
{
    name={\textbackslash neuespalte},
    description={
            Manueller Spaltenumbruch. Er ist nützlich wenn Kästen doof dargestellt werden oder Absätze in einer neuen Spalte beginnen sollensollen. Leitet direkt auf \textbackslash columnbreak weiter. \newline
            \textbf{Beispiel}: \textbackslash neuespalte
        }
}
\newglossaryentry{neueseite}
{
    name={\textbackslash neueseite},
    description={
            Manueller Seitenumbruch. Ist nützlich um, naja, den Beginn einer neuen Seite zu erzwingen..? Leitet direkt auf \textbackslash pagebreak weiter. \newline
            \textbf{Beispiel}: \textbackslash neueseite
        }
}

\newglossaryentry{tabelle}
{
    name={\textbackslash tabelle},
    description={\textbf{Beispiel}: \textbackslash tabelle\{X X X\}\{ eins \& zwei \& drei \textbackslash \textbackslash \}\{\}\newline
        }
}
\newglossaryentry{tkopf}
{
    name={\textbackslash tkopf},
    description={
            Setzt Schriftart und Farbe für einzelne Zellen in der ersten Reihe einer Tabelle. Muss auf jeden Eintrag (zwischen den \& einzeln angewendet werden)\newline
            \textbf{Beispiel}: siehe Tabelle im template.
        }
}
\newglossaryentry{kasten}
{
    name={\textbackslash kasten},
    description={
            Ein Pergamentkasten. In die Klammern können auch weitere Befehle eingesetzt werden, komplexe Umgebungen wie Tabellen könnten allerdings Probleme machen. Einfach Testen und Probleme berichten.\newline
            \textbf{Beispiel}: \textbackslash kasten\{Kastentext\}
        }
}
\newglossaryentry{kastengrau}
{
    name={\textbackslash kastengrau},
    description={
            Das selbe wie \textbackslash kasten nur mit blasserem Hintergrund.
            \newline
            \textbf{Beispiel}: \textbackslash kasten\{Kastentext\}
        }
}
\newglossaryentry{}
{
    name={\textbackslash },
    description={\textbf{Beispiel}: \textbackslash \{\}\newline
        }
}




%% befehle
\newcommand{\trennlinie}{\hr}
\newcommand{\anfang}{\raggedcolumns\frontmatter}
\newcommand{\hauptteil}{\mainmatter}
\newcommand{\schluss}{\backmatter}
\newcommand{\spaltenanfang}{\begin{multicols}{2}}
        \newcommand{\spaltenende}{\end{multicols}}
\newcommand{\skapitel}[1]{\end{multicols}\kapitel{#1}\begin{multicols}{2}}
\newcommand{\kapitel}[1]{{\let\cleardoublepage\clearpage\chapter*{#1}}\addcontentsline{toc}{chapter}{#1}}
\newcommand{\abschnitt}[1]{\section{#1}}
\newcommand{\leereseite}{\newpage\makebox{}\clearpage}
\newcommand{\platz}{\vfill}
\newcommand{\begriff}[1]{\textbf{#1}}
\newcommand{\geschichte}[2]{%
    \begin{center}
        \begin{minipage}[b]{0.9\linewidth}{\em #1}
            \begin{flushright}--- #2 \end{flushright}
        \end{minipage}
    \end{center}}
\newcommand{\kasten}[1]{\begin{pergamentbox}#1\end{pergamentbox}}
\newcommand{\inhaltsverzeichnis}{{\let\cleardoublepage\clearpage\tableofcontents}}
\newcommand{\fanprodukt}{\begin{center}\includegraphics[width=5cm]{layout/dsa_logo_fanprodukt}\end{center}}
\newcommand{\titelbild}[1]{
    \tikz[remember picture,overlay] \node[inner sep=0pt] at (current page.center){\includegraphics[width=\paperwidth,height=\paperheight]{#1}};}
\newcommand{\titel}[1]{
    \title{#1}
    \begin{centering}
        % 26.12.2020 \\
        \color{orange!15}
        %\vspace{0.55\textheight}
        \Huge
        %\fontfamily{\aniron}
        \fontsize{42}{42}
        \selectfont
        \aniron
        \shadowcolor{black}
        \shadowtext{#1} \\
    \end{centering}
}
\newcommand{\authorin}[1]{\begin{centering}\author{#1}\Large #1\\\end{centering}}
\newcommand{\titelseite}[1]{\thispagestyle{empty}#1}
%\newcommand{\kreatur}[1]{\begin{creature}#1\end{creature}}


\newcommand{\kreaturinfo}[2]{
    \footnotesize\textbf{#1: } #2}

\newcommand{\kreatur}[4]{%
    \begin{creaturebox}
        \begin{minipage}[b]{0.86\linewidth}
            {\color{dunkelrot}\Large\textbf{#1}}\\[1mm]
            \footnotesize \emph{#2}
        \end{minipage}
        \hfill\includegraphics[width=0.1\linewidth]{#3}
        \hr[0.8pt]
        \tcblower
        \vspace{2mm}
        %\tcbline
        #4
        \hr
    \end{creaturebox}
}

\newcommand{\credits}[1]{%
    \vfill
    \begin{center}
        #1
        % DAS SCHWARZE AUGE, AVENTURIEN und DERE sind eingetragene Marken... TODO: Copyright angabe?
    \end{center}
}
\newcommand{\absatz}[1]{\subsection*{#1}}
\newcommand{\spaltenumbruch}{\columnbreak}

\newcommand{\kreaturbaer}{\kreatur{Bär}{Raubtier der Wälder; großer Gegner}{gfx/kreaturen/tier}{
        \kreaturkampfwerte{11}{6}{6}{2}\trennlinie\kreaturwaffe{Prankenhieb}{1}{6}{13}{}{}\kreaturwaffe{Biss}{0}{3}{14}{4W6+6}{}\kreaturkampfvorteile{Standfest, Sturmangriff}\trennlinie\kreaturattribute{GE 16, KK 40, KO 42, MU 16}
    }}
\newcommand{\kreaturbaumdrache}{\kreatur{Baumdrache}{tierhafter Drache ohne magische Fähigkeiten}{gfx/kreaturen/tier}{
        \kreaturkampfwerte{4}{8}{1}{5}\trennlinie\kreaturvorteile{Flugfähig, Resistenz I (Feuer)}\trennlinie\kreaturwaffe{Klauen}{0}{6}{10}{2W6+1}{}\kreaturwaffe{Feueratem}{}{}{}{}{}\kreaturkampfvorteile{Niederwerfen, Sturmangriff (max. +6)}\trennlinie\kreaturattribute{GE 20, KK 6, KO 6, MU 8}
    }}
\newcommand{\kreaturblutfischschwarm}{\kreatur{Blutfischschwarm}{gierige Fressmaschinen der Dschungelflüsse; großer Gegner}{gfx/kreaturen/tier}{
        \kreaturkampfwerte{3}{2}{0}{1}\trennlinie\kreaturvorteile{Wasserwesen, Schmerzimmun II, Resistenz II (Stichwaffen), Verwundbarkeit IV (Flächenschaden)}\trennlinie\kreaturwaffe{Biss}{0}{4}{14}{2W6+3}{}\kreaturkampfvorteile{Zusätzliche Attacke II}
    }}
\newcommand{\kreaturelefant}{\kreatur{Elefant}{gutmütiger Pflanzenfresser; sehr großer Gegner}{gfx/kreaturen/tier}{
        \kreaturkampfwerte{15}{4}{9}{1}\trennlinie\kreaturwaffe{Stoßzähne}{1}{6}{8}{5W6+4}{}\kreaturwaffe{Trampeln}{1}{4}{6}{3W20}{}\kreaturkampfvorteile{Standfest, Überrennen, Unaufhaltsam, Zerstörerisch I und II}\trennlinie\kreaturattribute{KK 120, KO 140, MU 10}
    }}
\newcommand{\kreaturelefantelefantenruestung}{\kreatur{Elefant mit Rüstung}{gutmütiger Pflanzenfresser; sehr großer Gegner}{gfx/kreaturen/tier}{
        \kreaturkampfwerte{15}{4}{-1}{1}\trennlinie\kreaturwaffe{Stoßzähne}{1}{6}{8}{5W6+4}{}\kreaturwaffe{Trampeln}{1}{4}{6}{3W20}{}\kreaturkampfvorteile{Standfest, Überrennen, Unaufhaltsam, Zerstörerisch I und II}\trennlinie\kreaturattribute{KK 120, KO 140, MU 10}
    }}
\newcommand{\kreaturelefantkriegselefant}{\kreatur{Kriegselefant}{gutmütiger Pflanzenfresser; sehr großer Gegner}{gfx/kreaturen/tier}{
        \kreaturkampfwerte{15}{4}{9}{1}\trennlinie\kreaturwaffe{Stoßzähne}{1}{10}{12}{5W6+4}{}\kreaturwaffe{Trampeln}{1}{8}{10}{3W20}{}\kreaturkampfvorteile{Standfest, Überrennen, Unaufhaltsam, Zerstörerisch I und II}\trennlinie\kreaturattribute{KK 120, KO 140, MU 10}
    }}
\newcommand{\kreaturelefantmamut}{\kreatur{Mamut}{gutmütiger Pflanzenfresser; sehr großer Gegner}{gfx/kreaturen/tier}{
        \kreaturkampfwerte{17}{4}{9}{1}\trennlinie\kreaturwaffe{Stoßzähne}{1}{6}{8}{5W6+4}{}\kreaturwaffe{Trampeln}{1}{4}{6}{3W20}{}\kreaturkampfvorteile{Standfest, Überrennen, Unaufhaltsam, Zerstörerisch I und II, Resistenz gegen Kälte}\trennlinie\kreaturattribute{KK 160, KO 200, MU 10}
    }}
\newcommand{\kreaturfernhaendlerpaktierer}{\kreatur{Fernhändler und Paktierer}{ein gewitzter Händler und Paktierer der Charyptoroth}{gfx/kreaturen/tier}{
        \kreaturkampfwerte{5}{9}{4}{5}\trennlinie\kreaturvorteile{Gefahreninstinkt, Kreis der Verdammnis I, Willensstark I}\trennlinie\kreaturwaffe{Dolch}{0}{11}{11}{1W6+1}{}\kreaturkampfvorteile{Durchatmen}\trennlinie\kreaturattribute{CH 14, FF 88, GE 6, IN 12, KK 4, KL 8, KO 8, MU 12}
    }}
\newcommand{\kreaturghul}{\kreatur{Ghul}{ansteckender und gieriger Leichenfresser}{gfx/kreaturen/tier}{
        \kreaturkampfwerte{4}{9}{6}{6}\trennlinie\kreaturvorteile{Empfindlichkeit II (Sonnenlicht), Immunität (Einfluss), Lichtscheu, Schreckgestalt I}\trennlinie\kreaturwaffe{Unbewaffnet}{1}{6}{6}{1W6}{}\kreaturwaffe{Zähne}{0}{6}{6}{2W6+0}{}\trennlinie\kreaturattribute{GE 8, KK 8, KL 2, KO 8, MU 8}
    }}
\newcommand{\kreaturgoblin}{\kreatur{Goblin}{rotpelzige Landplage des Nordens}{gfx/kreaturen/tier}{
        \kreaturkampfwerte{4}{4}{5}{3}\trennlinie\kreaturvorteile{Angepasst I (Dunkelheit), Resistenz gegen Kälte}\trennlinie\kreaturwaffe{Holzspeer}{2}{8}{8}{2W6-1}{}\kreaturwaffe{Keule}{1}{8}{8}{2W6}{}\kreaturwaffe{Kurzbogen}{16}{}{}{2W6+1}{}\trennlinie\kreaturattribute{CH 4, FF 10, GE 8, IN 6, KK 4, KL 4, KO 4, MU 4}
    }}
\newcommand{\kreaturgreif}{\kreatur{Greif}{Mythenwesen und mächtiger Diener Praios‘; großer Gegner}{gfx/kreaturen/tier}{
        \kreaturkampfwerte{14}{20}{6}{8}\trennlinie\kreaturvorteile{Flugfähig, Willensstark II, Geisterpanzer (WS 20), Unbeugsamkeit}\trennlinie\kreaturwaffe{Prankenhieb}{1}{11}{17}{2W6+4}{}\kreaturwaffe{Schnabelhieb}{0}{2}{17}{3W6+4}{}\kreaturkampfvorteile{Zusätzlicher Attacke I, Haltet stand!, Hammerschlag, Niederwerfen, Standfest, Sturmangriff}\trennlinie\kreaturattribute{GE 16, KK 24, KO 26, MU 40}
    }}
\newcommand{\kreaturharpyie}{\kreatur{Harpyie}{wahnsinnige Verschmelzung von Vogel und Frau}{gfx/kreaturen/tier}{
        \kreaturkampfwerte{4}{4}{2}{8}\trennlinie\kreaturvorteile{Flieger}\trennlinie\kreaturwaffe{Klauen}{0}{8}{12}{2W6}{}\kreaturkampfvorteile{Niederwerfen, Sturmangriff}\trennlinie\kreaturattribute{GE 16, KK 6, KO 4, MU 4}
    }}
\newcommand{\kreaturhoehlenspinne}{\kreatur{Höhlenspinne}{einzelgängerische Spinne mit Lähmungsgift}{gfx/kreaturen/tier}{
        \kreaturkampfwerte{3}{9}{3}{3}\trennlinie\kreaturvorteile{Angepasst II (Dunkelheit), Immunität (Niederwerfen, Umreißen und ähnliche Effekte)}\trennlinie\kreaturwaffe{Biss}{0}{6}{9}{1W6+2}{}\trennlinie\kreaturattribute{KK 4, KO 4, MU 4}
    }}
\newcommand{\kreaturhornechse}{\kreatur{Hornechse}{friedfertiger Pflanzenfresser und Reittier der Achaz; großer Gegner}{gfx/kreaturen/tier}{
        \kreaturvorteile{Angepasst I (Sumpf), Verwundbarkeit I (Eis)}\trennlinie\kreaturwaffe{Horn}{0}{4}{8}{4W6+6}{}\kreaturkampfvorteile{Sturmangriff, Überrennen}\trennlinie\kreaturattribute{GE 6, KK 48, KO 54, MU 12}
    }}
\newcommand{\kreaturhyralkor}{\kreatur{Hyralkor, ein Kaiserdrache}{Herrscher seines Landstrichs; sehr großer Endgegner}{gfx/kreaturen/tier}{
        \kreaturkampfwerte{-1}{17}{4}{9}\trennlinie\kreaturvorteile{Aura (Hitze, Zähigkeit (28) jede INI:phase, 1 Wunde), Flugfähig, Immunität (Feuer), Schreckgestalt II, Verwundbarkeit I (Eis)}\trennlinie\kreaturwaffe{Biss}{2}{10}{19}{2W20+10}{}\kreaturwaffe{Flammenstrahl}{32}{}{}{5W6+8}{}\kreaturwaffe{Flammeninferno}{16}{}{}{4W6+4}{}\kreaturwaffe{Trampeln (Beine)}{0}{6}{11}{3W20}{}\kreaturwaffe{Pranken (Beine)}{1}{10}{16}{2W20+5}{}\kreaturwaffe{Schwanzhieb (S.)}{8}{8}{16}{4W6+2}{}\kreaturwaffe{Windstoß (Flügel)}{4}{6}{20}{}{}\kreaturkampfvorteile{Hammerschlag, Offensiver Kampfstil, Sturmangriff, Zusätzliche Attacke I}\trennlinie\kreaturattribute{GE 12, KK 200, KL 20, KO 220, MU 30}\trennlinie\keraturinfo{vorgehen}{Hyralkor bleibt auf Distanz und nutzt seine Flammenangriffe. Nur wenn seine Flügel schwer verletzt wurden oder er in beengten Verhältnissen kämpft, geht Hyralkor in den Nahkampf über – bevorzugt mit einem Sturmangriff aus der Luft. Am Boden verwendet Hyralkor seine Feuerangriffe gegen entfernte Gegner und seinen Prankenhieb gegen Angreifer in Schlagreichweite. Kämpfer unter sich versucht er zu zertrampeln, während er Angreifer hinter sich mit Schwanzhieben zur Seite fegt. Rücken ihm zu viele Angreifer zu nahe, versucht er sie mit einem Windstoß umzuwerfen um sich dann mit einem gewaltigen Satz Luft zu verschaffen. Sein Gebiss setzt er nur gegen schwer angeschlagene oder sehr widerstandsfähige Gegner ein, da er schmerzhafte Gegenangriffe in den Rachen befürchtet. Schwachpunkte: Um den Drachen zu töten, muss sein Rumpf oder Kopf vernichtet werden. Ersterer ist extrem zäh, während zweiterer für Nahkampfwaffen unerreichbar ist. Deswegen sollten zuerst anderer Zonen angegriffen werden: Ohne die Flügel wird Hyralkor auf den Boden gezwungen und wenn seine Beine einknicken, kommt der Kopf in Reichweite. Dabei hilft eine alte Verletzung Hyralkors: Einst rammten ihm die Brüder Turgosch und Targosch einen Drachentöter durch den rechten Hinterlauf. Die Spitze der Waffe steckt noch immer in der Wunde. Kennt ein Charakter die zwergische Geschichte, kann er mit einer GE- oder Akrobatik-Probe (24) in Position gelangen und die Waffe mit einer AT:  –12 weiter ins Bein treiben. Dadurch wird der Drache für eine INI:phase handlungsunfähig und heißes Blut schießt aus der Wunde (Akrobatik (24) oder 4W6 Schadenspunkte). Kampfunfähigkeit: Hyralkor flieht, wenn er 6 Wunden erlitten hat. In seinem Hort kämpft er jedoch bis zum Tod.}
    }}
\newcommand{\kreaturkrakenmolch}{\kreatur{Krakenmolch}{vielarmiger Schrecken der Küsten und Sümpfe; großer Gegner}{gfx/kreaturen/tier}{
        \kreaturkampfwerte{-1}{16}{2}{2}\trennlinie\kreaturvorteile{Amphibisch, Angepasst II (Wasser), Schreckgestalt I}\trennlinie\kreaturwaffe{Biss}{1}{4}{12}{4W6+4}{}\kreaturwaffe{Tentakelhieb}{8}{10}{14}{2W6+4}{}\kreaturwaffe{Tentakelklammer}{8}{10}{13}{2W6}{}\kreaturkampfvorteile{Zusätzliche Attacke II}\trennlinie\kreaturattribute{KK 12, KO 100, MU 20}
    }}
\newcommand{\kreaturkriegermumie}{\kreatur{Kriegermumie}{starker Wächter und letzter Anblick für zahlreiche Grabräuber}{gfx/kreaturen/tier}{
        \kreaturkampfwerte{10}{12}{4}{5}\trennlinie\kreaturvorteile{Astralsinn, Körperlosigkeit, Resistenz I (Stichwaffen), Verwundbarkeit II (Feuer), Schmerzimmun II, Schreckgestalt II}\trennlinie\kreaturwaffe{Khunchomer}{1}{10}{16}{5W6+4}{}\kreaturwaffe{Hände}{0}{10}{14}{2W6+6}{}\kreaturkampfvorteile{Kraftvoller Kampf III (Befreiungsschlag), Hammerschlag, Niederwerfen}\trennlinie\kreaturattribute{GE 12, KK 32}\trennlinie\keraturinfo{beispieldienste}{Kampf (1 Minute); Wache (1 Woche, mit Totes handle! permanent, –4)}\keraturinfo{beschwörung}{Skelettarius oder Totes handle! (20) aus einem mumifizierten Leichnam, 24 AsP}
    }}
\newcommand{\kreaturkriegspferd}{\kreatur{Kriegspferd}{kampferprobtes Reittier aus den Ställen Elenvinas; großer Gegner}{gfx/kreaturen/tier}{
        \kreaturkampfwerte{15}{6}{10}{1}\trennlinie\kreaturwaffe{Tritt}{1}{5}{12}{2W6+2}{}\kreaturwaffe{Biss}{0}{2}{12}{1W6+3}{}\kreaturkampfvorteile{Rüstungsgewöhnung, Standfest, im Reiterkampf WM +1}\trennlinie\kreaturattribute{GE 8, KK 26, KO 24, MU 12}
    }}
\newcommand{\kreaturkriegspferdleichtepferderuestung}{\kreatur{Kriegspferd mit leichter Rüstung}{kampferprobtes Reittier aus den Ställen Elenvinas; großer Gegner}{gfx/kreaturen/tier}{
        \kreaturkampfwerte{15}{6}{10}{1}\trennlinie\kreaturwaffe{Tritt}{1}{5}{12}{2W6+2}{}\kreaturwaffe{Biss}{0}{2}{12}{1W6+3}{}\kreaturkampfvorteile{Rüstungsgewöhnung, Standfest, im Reiterkampf WM +1}\trennlinie\kreaturattribute{GE 8, KK 26, KO 24, MU 12}
    }}
\newcommand{\kreaturkriegspferdmittlerepferderuestung}{\kreatur{Kriegspferd mit mittlerer Rüstung}{kampferprobtes Reittier aus den Ställen Elenvinas; großer Gegner}{gfx/kreaturen/tier}{
        \kreaturkampfwerte{15}{6}{9}{1}\trennlinie\kreaturwaffe{Tritt}{1}{4}{11}{2W6+2}{}\kreaturwaffe{Biss}{0}{1}{11}{1W6+3}{}\kreaturkampfvorteile{Rüstungsgewöhnung, Standfest, im Reiterkampf WM +1}\trennlinie\kreaturattribute{GE 8, KK 26, KO 24, MU 12}
    }}
\newcommand{\kreaturkriegspferdreitpony}{\kreatur{Reitpony}{kampferprobtes Reittier aus den Ställen Elenvinas; großer Gegner}{gfx/kreaturen/tier}{
        \kreaturkampfwerte{11}{6}{10}{1}\trennlinie\kreaturwaffe{Tritt}{1}{5}{12}{2W6+2}{}\kreaturwaffe{Biss}{0}{2}{12}{1W6+3}{}\kreaturkampfvorteile{Rüstungsgewöhnung, Standfest, im Reiterkampf WM +1}\trennlinie\kreaturattribute{GE 8, KK 26, KO 24, MU 12}
    }}
\newcommand{\kreaturkriegspferdreittier}{\kreatur{Reitpferd}{kampferprobtes Reittier aus den Ställen Elenvinas; großer Gegner}{gfx/kreaturen/tier}{
        \kreaturkampfwerte{13}{6}{12}{1}\trennlinie\kreaturwaffe{Tritt}{1}{5}{12}{2W6+2}{}\kreaturwaffe{Biss}{0}{2}{12}{1W6+3}{}\kreaturkampfvorteile{Rüstungsgewöhnung, Standfest, im Reiterkampf WM +1}\trennlinie\kreaturattribute{GE 8, KK 26, KO 24, MU 12}
    }}
\newcommand{\kreaturkriegspferdtralloper}{\kreatur{Tralloper}{kampferprobtes Reittier aus den Ställen Elenvinas; großer Gegner}{gfx/kreaturen/tier}{
        \kreaturkampfwerte{16}{6}{8}{1}\trennlinie\kreaturwaffe{Tritt}{1}{5}{12}{2W6+2}{}\kreaturwaffe{Biss}{0}{2}{12}{1W6+3}{}\kreaturkampfvorteile{Rüstungsgewöhnung, Standfest, im Reiterkampf WM +1}\trennlinie\kreaturattribute{GE 8, KK 26, KO 24, MU 12}
    }}
\newcommand{\kreaturleibwaechter}{\kreatur{Leibwächter}{persönlicher Wächter eines Würdenträgers}{gfx/kreaturen/tier}{
        \kreaturkampfwerte{4}{5}{6}{12}\trennlinie\kreaturwaffe{Amazonensäbel}{1}{16}{13}{2W6+3}{}\kreaturwaffe{Holzschild}{0}{18}{15}{1W6+2}{}\kreaturwaffe{Faust}{0}{14}{14}{1W6+4}{}\kreaturkampfvorteile{Schildkampf III (Schildwall), Aufmerksamkeit, Ausfall, Defensiver Kampfstil, Kampfreflexe, Niederwerfen, Standfest, Sturmangriff, Waffenloser Kampf}\trennlinie\kreaturattribute{CH 6, FF 10, GE 12, IN 16, KK 16, KL 6, KO 10, MU 8}
    }}
\newcommand{\kreaturmorfu}{\kreatur{Morfu}{schleimige Giftschnecke, die ihre Stacheln verschießt; großer Gegner}{gfx/kreaturen/tier}{
        \kreaturkampfwerte{4}{16}{1}{0}\trennlinie\kreaturvorteile{Immunität (Niederwerfen, Umreißen und ähnliche Effekte), Resistenz I (Stumpfe Waffen)}\trennlinie\kreaturwaffe{Dornensalve}{8}{}{}{2W6+2}{}\trennlinie\kreaturattribute{KK 8, KO 8, MU 10}
    }}
\newcommand{\kreaturoger}{\kreatur{Oger}{stinkender Vielfraß und Menschenfresser; großer Gegner}{gfx/kreaturen/tier}{
        \kreaturkampfwerte{12}{8}{6}{4}\trennlinie\kreaturwaffe{Faust}{1}{8}{14}{2W6+4}{}\kreaturwaffe{Keule}{2}{8}{14}{4W6+6}{}\kreaturkampfvorteile{Kraftvoller Kampf III (Befreiungsschlag), Niederwerfen, Sturmangriff}\trennlinie\kreaturattribute{GE 4, KK 22, KL –2, KO 20, MU 20}
    }}
\newcommand{\kreaturogerkriegsoger}{\kreatur{Kriegsoger}{stinkender Vielfraß und Menschenfresser; großer Gegner}{gfx/kreaturen/tier}{
        \kreaturkampfwerte{14}{10}{8}{6}\trennlinie\kreaturwaffe{Faust}{1}{8}{14}{2W6+4}{}\kreaturwaffe{Keule}{2}{8}{14}{4W6+6}{}\kreaturkampfvorteile{Kraftvoller Kampf III (Befreiungsschlag), Niederwerfen, Sturmangriff}\trennlinie\kreaturattribute{GE 4, KK 22, KL –2, KO 20, MU 20}
    }}
\newcommand{\kreaturogerstatue}{\kreatur{Ogerstatue}{behäbige Bronzestatue eines Ogers; großer Gegner}{gfx/kreaturen/tier}{
        \kreaturkampfwerte{10}{20}{1}{-1}\trennlinie\kreaturvorteile{Immunität (Gifte, Krankheiten), Resistenz II (Feuer, Stichwaffen)}\trennlinie\kreaturwaffe{Faust}{2}{4}{12}{2W20+6}{}\kreaturkampfvorteile{Zusätzliche Attacke I, Unaufhaltsam}\trennlinie\kreaturattribute{KK 44, KO 48}\trennlinie\keraturinfo{beispieldienste}{Lastentransport (+4, 1 Tag); Wache (1 Tag); Kampf (–4, 1 Minute)}\keraturinfo{beschwörung}{Stein wandle! (28) aus Bronzeplatten, 32 AsP}
    }}
\newcommand{\kreaturorkkrieger}{\kreatur{Ork-Krieger}{durchschnittlicher orkischer Kämpfer}{gfx/kreaturen/tier}{
        \kreaturkampfwerte{5}{5}{4}{3}\trennlinie\kreaturwaffe{Arbach}{1}{10}{10}{2W6+6}{}\kreaturwaffe{Kriegshammer}{1}{10}{10}{3W6+6}{}\kreaturwaffe{Faust}{0}{9}{9}{1W6+2}{}\kreaturkampfvorteile{Kraftvoller Kampf II, Natürliche Rüstung, Niederwerfen, Offensiver Kampfstil, Rüstungsgewöhnung, Waffenloser Kampf}\trennlinie\kreaturattribute{CH 2, FF 4, GE 8, IN 6, KK 12, KL 2, KO 12, MU 14}
    }}
\newcommand{\kreaturorkkundschafter}{\kreatur{Ork-Kundschafter}{wildniserfahrener orkischer Kundschafter oder Räuber}{gfx/kreaturen/tier}{
        \kreaturkampfwerte{5}{5}{5}{4}\trennlinie\kreaturwaffe{Skraja}{0}{8}{8}{2W6+2}{}\kreaturwaffe{Ork. Reiterbogen}{32}{}{}{2W6+3}{}\kreaturkampfvorteile{Natürliche Rüstung, Scharfschuss, Reflexschuss}\trennlinie\kreaturattribute{CH 2, FF 12, GE 8, IN 8, KK 10, KL 2, KO 12, MU 10}
    }}
\newcommand{\kreaturperldrache}{\kreatur{Perldrache}{intelligenter Drache und geschickter Fischer; großer Gegner}{gfx/kreaturen/tier}{
        \kreaturkampfwerte{8}{12}{2}{6}\trennlinie\kreaturvorteile{Angepasst I (Wasser), Aura (Hitze, Zähigkeit (20) alle 4 INI:phasen, 1 Wunde), Flugfähig, Resistenz II (Feuer), Schreckgestalt I}\trennlinie\kreaturwaffe{Biss}{2}{8}{12}{5W6+4}{}\kreaturwaffe{Klauenhieb}{1}{12}{12}{4W6+4}{}\kreaturwaffe{Schwanzschlag}{2}{8}{14}{2W6+2}{}\kreaturkampfvorteile{Niederwerfen, Standfest, Sturmangriff, Zerstörerisch I, Zusätzliche Attacke I}\trennlinie\kreaturattribute{GE 12, KK 70, KO 74, MU 20}
    }}
\newcommand{\kreaturperldrachewestwinddrache}{\kreatur{Westwinddrache}{intelligenter Drache und geschickter Fischer; großer Gegner}{gfx/kreaturen/tier}{
        \kreaturkampfwerte{8}{12}{2}{6}\trennlinie\kreaturvorteile{Angepasst I (Wasser), Aura (Hitze, Zähigkeit (20) alle 4 INI:phasen, 1 Wunde), Flugfähig, Resistenz II (Feuer), Schreckgestalt I}\trennlinie\kreaturwaffe{Biss}{2}{8}{12}{5W6+4}{}\kreaturwaffe{Klauenhieb}{1}{12}{12}{4W6+4}{}\kreaturwaffe{Schwanzschlag}{2}{8}{14}{2W6+2}{}\kreaturkampfvorteile{Niederwerfen, Standfest, Sturmangriff, Zerstörerisch I, Zusätzliche Attacke I}\trennlinie\kreaturattribute{GE 12, KK 70, KO 74, MU 20}
    }}
\newcommand{\kreaturpestrattenschwarm}{\kreatur{Pestrattenschwarm}{ansteckender Schwarm aus Rattenkadavern; großer Gegner}{gfx/kreaturen/tier}{
        \kreaturkampfwerte{4}{2}{2}{0}\trennlinie\kreaturvorteile{Astralsinn, Resistenz II (Stichwaffen), Schmerzimmun II, Verwundbarkeit IV (Flächenschaden)}\trennlinie\kreaturwaffe{Biss}{0}{6}{14}{2W6–1}{}\kreaturkampfvorteile{Zusätzliche Attacke I}\trennlinie\keraturinfo{beispieldienste}{Stall leerfessen (1 Stunde, +4); andere Lebewesen fressen (1 Minute); nur bestimmtes Lebewesen fressen (1 Minute, –4)}\keraturinfo{beschwörung}{Skelettarius oder Totes handle! (16) aus Rattenkadavern, 16 AsP}
    }}
\newcommand{\kreaturpestrattenschwarmsehrgrosserschwarm}{\kreatur{Sehr großer Pestrattenschwarm}{ansteckender Schwarm aus Rattenkadavern; großer Gegner}{gfx/kreaturen/tier}{
        \kreaturkampfwerte{5}{2}{2}{0}\trennlinie\kreaturvorteile{Astralsinn, Resistenz II (Stichwaffen), Schmerzimmun II, Verwundbarkeit IV (Flächenschaden)}\trennlinie\kreaturwaffe{Biss}{0}{6}{16}{2W6–1}{}\kreaturkampfvorteile{Zusätzliche Attacke II}\trennlinie\keraturinfo{beispieldienste}{Stall leerfessen (1 Stunde, +4); andere Lebewesen fressen (1 Minute); nur bestimmtes Lebewesen fressen (1 Minute, –4)}\keraturinfo{beschwörung}{Skelettarius oder Totes handle! (20) aus Rattenkadavern, 24 AsP}
    }}
\newcommand{\kreaturraeuber}{\kreatur{Räuber}{ausgehungerte Halsabschneider auf der Suche nach leichter Beute}{gfx/kreaturen/tier}{
        \kreaturkampfwerte{4}{5}{4}{2}\trennlinie\kreaturvorteile{Zweihändig}\trennlinie\kreaturwaffe{Keule}{1}{8}{8}{2W6}{}\kreaturwaffe{Holzspeer}{2}{8}{8}{2W6+1}{}\kreaturwaffe{Kurzbogen}{16}{}{}{2W6+1}{}\trennlinie\kreaturattribute{CH 4, FF 8, GE 6, IN 4, KK 6, KL 6, KO 6, MU 8}
    }}
\newcommand{\kreaturriese}{\kreatur{Riese}{Unsterblicher Gigant; sehr großer Gegner}{gfx/kreaturen/tier}{
        \kreaturkampfwerte{16}{16}{9}{6}\trennlinie\kreaturvorteile{Schreckgestalt II, Unbeugsamkeit, Willensstark II}\trennlinie\kreaturwaffe{Faust}{2}{12}{16}{2W20+5}{}\kreaturwaffe{Baumstamm}{5}{12}{18}{3W20+10}{}\kreaturwaffe{Trampeln}{1}{6}{16}{4W20}{}\kreaturwaffe{Felsbrocken}{16}{}{14}{3W20}{}\kreaturkampfvorteile{Zusätzliche Attacke I (nur Trampeln), Standfest, Sturmangriff, Unaufhaltsam, Zerstörerisch I und II}\trennlinie\kreaturattribute{GE 6, KK 150, KO 130, MU 20}\trennlinie\keraturinfo{schwachpunkte}{Im Nacken des Riesen, vom Boden unsichtbar, liegt eine armdicke Ader direkt unter der Haut. Sie kann entweder von einem Fernkämpfer in ähnlicher Höhe mit Sinnenschärfe (28) erspäht und mit einem FK: (28) getroffen werden, oder ein tollkühner Kletterer erklimmt den Riesen (Klettern (28)). Nach einem Treffer müssen sich nur noch alle Umstehenden in Sicherheit bringen, wenn sich der Riese in Schmerzen zu Boden wirft und alles niederwalzt (in 3 Schritt Umkreis GE: oder Akrobatik 30, sonst 3W20 TP:  ). Das kostet den Riesen 1 Aktion und gibt Umstehenden die Chance auf einen Passierschlag.}\keraturinfo{vorgehen}{Eine Übermacht hält sich der Riese mit Befreiungsschlägen mit einem Baumstamm vom Leib, während er die Gegner unter sich niedertrampelt. Wegen des Vorteils Unaufhaltsam sollte seinen Angriffen unbedingt ausgewichen werden. Sollte ihm ein Fernkämpfer oder Zauberer gefährlich werden, stürmt er auf ihn zu (und trampelt alles im Weg nieder) oder bewirft ihn mit gewaltigen Felsbrocken.}
    }}
\newcommand{\kreaturriesenamoebe}{\kreatur{Riesenamöbe}{allesfressender Schleimhaufen; großer Gegner}{gfx/kreaturen/tier}{
        \kreaturkampfwerte{6}{14}{1}{0}\trennlinie\kreaturvorteile{Immunität (Niederwerfen, Umreißen und ähnliche Effekte), Körperlosigkeit, Resistenz I (Stichwaffen, stumpfe Waffen), Schmerzimmun II}\trennlinie\kreaturwaffe{Scheinarm}{1}{2}{6}{1W6+0}{}
    }}
\newcommand{\kreaturriesenkaiman}{\kreatur{Riesenkaiman}{gefährlicher Jäger der südlichen Gewässer}{gfx/kreaturen/tier}{
        \kreaturkampfwerte{13}{4}{4}{4}\trennlinie\kreaturvorteile{Amphibisch, Immunität (Niederwerfen, Umreißen und ähnliche Effekte), Verwundbarkeit (Eis)}\trennlinie\kreaturwaffe{Biss}{0}{6}{8}{4W6+4}{}\kreaturwaffe{Schwanz}{2}{2}{6}{1W6+2}{}\kreaturkampfvorteile{Zusätzliche Attacke I}\trennlinie\kreaturattribute{KK 20, KO 38, MU 10}
    }}
\newcommand{\kreaturriesenkraken}{\kreatur{Riesenkraken}{Albtraum vieler Seefahrer; sehr großer Endgegner}{gfx/kreaturen/tier}{
        \kreaturkampfwerte{-1}{20}{6}{2}\trennlinie\kreaturvorteile{Schreckgestalt III, Wasserwesen, Willensstark II}\trennlinie\kreaturwaffe{Biss (Rumpf)}{1}{2}{14}{3W20+10}{}\kreaturwaffe{Tentakelhieb}{12}{8}{18}{4W6+2}{}\kreaturwaffe{Tentakelklammer}{10}{8}{15}{4W6}{}\kreaturkampfvorteile{Zusätzliche Attacke IV (nur Tentakel), Unaufhaltsam}\trennlinie\kreaturattribute{KK 22, KO 220}\trennlinie\keraturinfo{schwachpunkte}{Der Riesenkraken greift aus der Tiefe an, weswegen er hervorragend gegen Angriffen von oben geschützt ist. Seine Unterseite birgt jedoch ein verwundbares Verdauungsorgan, das allerdings erst erreicht werden muss (Schwimmen (20) und eine Aktion zum Hinabtauchen, Schwimmen (24) und eine weitere Aktion, um in Position zu kommen.) Auch eine Amphore Brandöl kann dem Kraken schwerste Verbrennungen zufügen, wenn sie direkt in das Maul geworfen wird. Für einen solchen Wurf ist eine IN-Probe (28) nötig, um die Bewegung des Mauls einzuschätzen. Der folgende Wurf erfordert eine KK-Probe (28 oder mehr). Die Verwendung von Telekinesemagie ersetzt die KK-Probe, nicht aber die IN-Probe. Die Verbrennungen führen zu einer Tobsucht, bei der der Kraken alle umklammerten Gegner loslässt und mit 6 Tentakeln und voller Offensive auf das Schiff einschlägt. Kampfunfähigkeit: Der Kraken flieht, wenn er 4 Wunden erlitten oder die Hälfte seiner Tentakel verloren hat.}\keraturinfo{vorgehen}{Mit ungefähr der Hälfte seiner Tentakel hält der Riesenkraken das Schiff fest, während er mit den anderen das Deck verheert. Dazu nutzt er mindestens drei seiner zusätzlichen Attacken als Tentakelhiebe. Unterdessen versucht er, einzelne Opfer zu fressen: Zuerst umfasst er sie mit einer Tentakelklammer, in der nächsten Aktion zieht er sie zu seinem Maul und in der übernächsten Aktion beißt er zu.}
    }}
\newcommand{\kreaturritter}{\kreatur{Ritter}{erfahrener Veteran vieler Gefechte gegen den Ork}{gfx/kreaturen/humanoid}{
        \kreaturkampfwerte{5}{5}{2}{8}\trennlinie\kreaturwaffe{Breitschwert}{1}{11}{11}{2W6+6}{}\kreaturwaffe{Großschild}{0}{13}{13}{2W6+1}{}\kreaturwaffe{Andergaster}{2}{9}{9}{3W6+7}{}\kreaturwaffe{Kriegslanze}{2}{0}{9}{3W6+4}{}\kreaturwaffe{Tralloper}{1}{15}{15}{2W6+6}{}\kreaturkampfvorteile{Reiterkampf III* (VT, AT, TP:  +3), abgesessen Schildkampf III* (VT:  +3) oder Kraftvoller Kampf III* (TP:  +3), Ausfall, Haltet stand!, Hammerschlag, Kampfreflexe, Muskelprotz, Niederwerfen, Offensiver Kampfstil, Rüstungsgewöhnung, Durchatmen}\trennlinie\kreaturattribute{CH 10, FF 6, GE 8, IN 8, KK 16, KL 8, KO 12, MU 14}
    }}
\newcommand{\kreatursaebelzahntiger}{\kreatur{Säbelzahntiger}{König der Dschungel und nördlichen Steppen; großer Gegner}{gfx/kreaturen/tier}{
        \kreaturkampfwerte{11}{6}{11}{10}\trennlinie\kreaturvorteile{Angepasst II (Dunkelheit), Blitzschnell (erleidet keine Passierschläge, wenn er sich aus dem Nahkampf zurückzieht)}\trennlinie\kreaturwaffe{Prankenhieb}{1}{12}{16}{2W6+4}{}\kreaturwaffe{Biss}{0}{2}{16}{5W6+2}{}\kreaturkampfvorteile{Standfest, Sturmangriff}\trennlinie\kreaturattribute{GE 20, KK 20, KO 22, MU 18}
    }}
\newcommand{\kreatursaebelzahntigerloewe}{\kreatur{Löwe}{König der Dschungel und nördlichen Steppen; großer Gegner}{gfx/kreaturen/tier}{
        \kreaturkampfwerte{10}{6}{11}{10}\trennlinie\kreaturvorteile{Angepasst II (Dunkelheit), Blitzschnell (erleidet keine Passierschläge, wenn er sich aus dem Nahkampf zurückzieht)}\trennlinie\kreaturwaffe{Prankenhieb}{1}{12}{16}{2W6+0}{}\kreaturwaffe{Biss}{0}{2}{16}{5W6-2}{}\kreaturkampfvorteile{Standfest, Sturmangriff}\trennlinie\kreaturattribute{GE 16, KK 16, KO 18, MU 14}
    }}
\newcommand{\kreatursaebelzahntigerpanther}{\kreatur{Panther}{König der Dschungel und nördlichen Steppen; großer Gegner}{gfx/kreaturen/tier}{
        \kreaturkampfwerte{11}{6}{11}{10}\trennlinie\kreaturvorteile{Angepasst II (Dunkelheit), Blitzschnell (erleidet keine Passierschläge, wenn er sich aus dem Nahkampf zurückzieht)}\trennlinie\kreaturwaffe{Prankenhieb}{1}{12}{16}{2W6+4}{}\kreaturwaffe{Biss}{0}{2}{16}{5W6+2}{}\kreaturkampfvorteile{Standfest, Sturmangriff}\trennlinie\kreaturattribute{GE 20, KK 20, KO 22, MU 18}
    }}
\newcommand{\kreaturschroeter}{\kreatur{Schröter}{Landplage oder exotisches Haustier}{gfx/kreaturen/tier}{
        \kreaturkampfwerte{4}{13}{1}{0}\trennlinie\kreaturwaffe{Zange}{1}{4}{9}{2W6+2}{}\kreaturwaffe{Biss}{0}{2}{12}{2W6}{}\kreaturkampfvorteile{Doppelangriff, Standfest}\trennlinie\kreaturattribute{GE 4, KK 8, KO 14, MU 4}
    }}
\newcommand{\kreaturskelett}{\kreatur{Skelett}{schwaches untotes Kanonenfutter}{gfx/kreaturen/tier}{
        \kreaturkampfwerte{4}{5}{3}{2}\trennlinie\kreaturvorteile{Astralsinn, Resistenz II (Stichwaffen), Rudel}\trennlinie\kreaturwaffe{Hände}{0}{5}{8}{1W6+2}{}\kreaturwaffe{Streitaxt}{1}{5}{8}{2W6+3}{}\kreaturwaffe{Stoßspeer}{2}{4}{7}{3W6+4}{}\trennlinie\kreaturattribute{GE 4, KK 8}\trennlinie\keraturinfo{beispieldienste}{Burggraben auffüllen (1 Stunde, +4); Kampf (1 Minute); Wache (1 Tag, mit Totes handle! permanent, –4)}\keraturinfo{beschwörung}{Skelettarius oder Totes handle! (12) aus einer skelettierten Leiche, 8 AsP}
    }}
\newcommand{\kreaturskelettruestung}{\kreatur{Skelett mit Rüstung}{schwaches untotes Kanonenfutter}{gfx/kreaturen/tier}{
        \kreaturkampfwerte{5}{5}{2}{1}\trennlinie\kreaturvorteile{Astralsinn, Resistenz II (Stichwaffen), Rudel}\trennlinie\kreaturwaffe{Hände}{0}{5}{8}{1W6+2}{}\kreaturwaffe{Streitaxt}{1}{5}{8}{2W6+3}{}\kreaturwaffe{Stoßspeer}{2}{4}{7}{3W6+4}{}\trennlinie\kreaturattribute{GE 4, KK 8}\trennlinie\keraturinfo{beispieldienste}{Burggraben auffüllen (1 Stunde, +4); Kampf (1 Minute); Wache (1 Tag, mit Totes handle! permanent, –4)}\keraturinfo{beschwörung}{Skelettarius oder Totes handle! (12) aus einer skelettierten Leiche, 8 AsP}
    }}
\newcommand{\kreatursoeldner}{\kreatur{Söldner}{verlässlicher Mietling oder undisziplinierter Halsabschneider}{gfx/kreaturen/tier}{
        \kreaturkampfwerte{5}{5}{5}{3}\trennlinie\kreaturwaffe{Säbel}{1}{11}{11}{2W6+3}{}\kreaturwaffe{Streitaxt}{1}{11}{11}{2W6+3}{}\kreaturwaffe{Holzschild}{0}{12}{12}{1W6}{}\kreaturwaffe{Speer}{2}{11}{11}{3W6}{}\kreaturwaffe{Hellebarde}{2}{11}{11}{2W6+3}{}\kreaturkampfvorteile{Schildkampf I* (VT:  +1), Kraftvoller Kampf I* (TP:  +1), Schneller Kampf I* (AT:  +1), Niederwerfen, Rüstungsgewöhnung}\trennlinie\kreaturattribute{CH 6, FF 6, GE 8, IN 6, KK 10, KL 4, KO 10, MU 10}
    }}
\newcommand{\kreatursoeldnerkommandant}{\kreatur{Söldnerkommandant}{charismatischer Befehlshaber einer größeren Einheit}{gfx/kreaturen/tier}{
        \kreaturkampfwerte{5}{13}{5}{10}\trennlinie\kreaturvorteile{Willensstark I und II}\trennlinie\kreaturwaffe{Tuzakmesser}{2}{14}{16}{2W6+5}{}\kreaturwaffe{Kurzschwert}{0}{14}{16}{2W6}{}\kreaturwaffe{Faust}{0}{8}{8}{1W6+1}{}\kreaturkampfvorteile{Schneller Kampf II, Defensiver Kampfstil, Formiert euch!, Haltet stand!, Kampfreflexe, Keine Gefangenen!, Niederwerfen, Rüstungsgewöhnung, Standfest, Sturmangriff, Durchatmen}\trennlinie\kreaturattribute{CH 16, FF 6, GE 12, IN 12, KK 8, KL 8, KO 12, MU 14}
    }}
\newcommand{\kreatursoeldnerveteran}{\kreatur{Söldnerveran}{verlässlicher Mietling oder undisziplinierter Halsabschneider}{gfx/kreaturen/tier}{
        \kreaturkampfwerte{5}{5}{5}{8}\trennlinie\kreaturwaffe{Säbel}{1}{12}{13}{2W6+3}{}\kreaturwaffe{Streitaxt}{1}{12}{13}{2W6+3}{}\kreaturwaffe{Holzschild}{0}{13}{14}{1W6}{}\kreaturwaffe{Speer}{2}{12}{13}{3W6}{}\kreaturwaffe{Hellebarde}{2}{11}{13}{2W6+3}{}\kreaturkampfvorteile{Schildkampf I* (VT:  +1), Kraftvoller Kampf II* (TP:  +1), Schneller Kampf II* (AT:  +1), Niederwerfen, Rüstungsgewöhnung, Kampfreflexe, Standfest, Waffenloser Kampf, Durchatmen}\trennlinie\kreaturattribute{CH 8, FF 7, GE 8, IN 7, KK 11, KL 6, KO 10, MU 11}
    }}
\newcommand{\kreaturstadtwache}{\kreatur{Stadtwache}{stets beschäftigter Hüter von Recht und Ordnung}{gfx/kreaturen/tier}{
        \kreaturkampfwerte{4}{5}{4}{6}\trennlinie\kreaturwaffe{Hellebarde}{2}{7}{8}{2W6+2}{}\kreaturwaffe{Kurzschwert}{0}{9}{10}{2W6}{}\kreaturwaffe{Leichte Armbrust}{20}{}{}{3W6+1}{}\kreaturkampfvorteile{Schneller Kampf I, Standfest}\trennlinie\kreaturattribute{CH 8, FF 6, GE 12, IN 12, KK 6, KL 8, KO 6, MU 8}
    }}
\newcommand{\kreatursumupriester}{\kreatur{Sumupriester}{Beschützer der Natur, weiser Ratgeber, unerbittlicher Feind}{gfx/kreaturen/tier}{
        \kreaturkampfwerte{5}{5}{4}{8}\trennlinie\kreaturwaffe{Kampfstab}{2}{9}{9}{1W6+1}{}\kreaturwaffe{Vulkanglasdolch}{0}{4}{4}{1W6+1}{}\trennlinie\kreaturattribute{CH 8, FF 4, GE 4, IN 16, KK 4, KL 16, KO 8, MU 10}
    }}
\newcommand{\kreaturtatzelwurm}{\kreatur{Tatzelwurm}{stinkender Vertreter der niederen Drachen; großer Gegner}{gfx/kreaturen/tier}{
        \kreaturkampfwerte{10}{8}{3}{2}\trennlinie\kreaturvorteile{Aura (Gestank), Zähigkeit (16) jede INI:phase, –2 auf körperliche Proben (kumulativ) für 1 Tag), Resistenz I (Feuer), Schreckgestalt I}\trennlinie\kreaturwaffe{Biss}{1}{4}{15}{4W6+2}{}\kreaturwaffe{Klauen}{1}{6}{12}{3W6+1}{}\kreaturwaffe{Schwanz}{2}{4}{10}{2W6+1}{}\kreaturkampfvorteile{Zusätzliche Attacke I, Standfest}\trennlinie\kreaturattribute{KK 26, KO 32}
    }}
\newcommand{\kreaturtroll}{\kreatur{Troll}{jähzorniger Überlebender einer uralten Kultur; großer Gegner}{gfx/kreaturen/tier}{
        \kreaturkampfwerte{12}{8}{7}{3}\trennlinie\kreaturwaffe{Faust}{1}{9}{15}{4W6+1}{}\kreaturwaffe{Keule}{2}{8}{14}{5W6+4}{}\kreaturwaffe{Axt}{2}{9}{15}{6W6+4}{}\kreaturkampfvorteile{Ausfall, Kraftvoller Kampf III (Befreiungsschlag), Atemtechnik, Sturmangriff}\trennlinie\kreaturattribute{KK 32, KL 6, KO 34, MU 12}
    }}
\newcommand{\kreaturwaldschrat}{\kreatur{Waldschrat}{baumartiger Hüter der Wälder; großer Gegner}{gfx/kreaturen/tier}{
        \kreaturkampfwerte{10}{10}{6}{3}\trennlinie\kreaturwaffe{Asthieb}{1}{6}{15}{5W6}{}\kreaturkampfvorteile{Kraftvoller Kampf III (Befreiungsschlag), Atemtechnik, Standfest}\trennlinie\kreaturattribute{KK 32, KO 34, MU 12}
    }}
\newcommand{\kreaturwaldschratnichtmagisch}{\kreatur{Nichtmagischer Waldschrat}{baumartiger Hüter der Wälder; großer Gegner}{gfx/kreaturen/tier}{
        \kreaturkampfwerte{10}{10}{6}{3}\trennlinie\kreaturvorteile{Tarnung (im Wald), Verwundbarkeit I (Feuer)}\trennlinie\kreaturwaffe{Asthieb}{1}{6}{15}{5W6}{}\kreaturkampfvorteile{Kraftvoller Kampf III (Befreiungsschlag), Atemtechnik, Standfest}\trennlinie\kreaturattribute{KK 32, KO 34, MU 12}
    }}
\newcommand{\kreaturweisserhetzer}{\kreatur{Weisser Hetzer}{zweiköpfige Chimäre aus Wolf und Dämon}{gfx/kreaturen/tier}{
        \kreaturkampfwerte{4}{7}{8}{5}\trennlinie\kreaturvorteile{Angepasst II (Dunkelheit)}\trennlinie\kreaturwaffe{Biss}{1}{9}{11}{2W6+2}{}\kreaturkampfvorteile{Zusätzliche Attacke I, Niederwerfen, Standfest, Sturmangriff}\trennlinie\kreaturattribute{GE 16, KK 8, KO 4}\trennlinie\keraturinfo{beispieldienste}{Spur verfolgen (1 Stunde, +4); Kampf (1 Minute); Hetzen und Töten (1 Stunde, –4)}\keraturinfo{beschwörung}{Chimaeroform (16) aus einem Wolf und einem Karmanath oder Tuur-Amash, 16 AsP}
    }}
\newcommand{\kreaturwildschwein}{\kreatur{Wildschwein}{gefräßiger Schädling und goblinisches Reittier}{gfx/kreaturen/tier}{
        \kreaturkampfwerte{9}{3}{7}{3}\trennlinie\kreaturwaffe{Stoß}{0}{3}{10}{2W6+4}{}\kreaturkampfvorteile{Standfest, Sturmangriff, im Reiterkampf WM +0}\trennlinie\kreaturattribute{GE 8, KK 18, KO 22}
    }}
\newcommand{\kreaturwildschweinkriegswildschwein}{\kreatur{Kriegswildschwein}{gefräßiger Schädling und goblinisches Reittier}{gfx/kreaturen/tier}{
        \kreaturkampfwerte{11}{3}{7}{3}\trennlinie\kreaturwaffe{Stoß}{0}{4}{11}{2W6+6}{}\kreaturkampfvorteile{Standfest, Sturmangriff, im Reiterkampf WM +0}\trennlinie\kreaturattribute{GE 8, KK 18, KO 22}
    }}
\newcommand{\kreaturwolf}{\kreatur{Wolf}{verbreitetes Rudeltier}{gfx/kreaturen/tier}{
        \kreaturkampfwerte{4}{2}{8}{2}\trennlinie\kreaturvorteile{Angepasst II (Dunkelheit)}\trennlinie\kreaturwaffe{Biss}{0}{6}{6}{2W6}{}\kreaturkampfvorteile{Niederwerfen, Standfest, Sturmangriff}\trennlinie\kreaturattribute{GE 14, KK 6, KO 6, MU 6}
    }}
\newcommand{\kreaturwolfgrosserhund}{\kreatur{Großer Hund}{verbreitetes Rudeltier}{gfx/kreaturen/tier}{
        \kreaturkampfwerte{4}{2}{8}{2}\trennlinie\kreaturvorteile{Angepasst II (Dunkelheit)}\trennlinie\kreaturwaffe{Biss}{0}{6}{6}{2W6}{}\kreaturkampfvorteile{Niederwerfen, Standfest, Sturmangriff}\trennlinie\kreaturattribute{GE 14, KK 6, KO 6, MU 6}
    }}
\newcommand{\kreaturwolfkampfhund}{\kreatur{Kampfhund}{verbreitetes Rudeltier}{gfx/kreaturen/tier}{
        \kreaturkampfwerte{5}{3}{7}{3}\trennlinie\kreaturvorteile{Angepasst II (Dunkelheit)}\trennlinie\kreaturwaffe{Biss}{0}{6}{6}{2W6}{}\kreaturkampfvorteile{Niederwerfen, Standfest, Sturmangriff}\trennlinie\kreaturattribute{GE 14, KK 8, KO 8, MU 6}
    }}
\newcommand{\kreaturwolfkhoramsbestie}{\kreatur{Khoramsbestie}{verbreitetes Rudeltier}{gfx/kreaturen/tier}{
        \kreaturkampfwerte{4}{4}{10}{4}\trennlinie\kreaturvorteile{Angepasst II (Dunkelheit)}\trennlinie\kreaturwaffe{Biss}{0}{6}{6}{2W6}{}\kreaturkampfvorteile{Niederwerfen, Standfest, Sturmangriff}\trennlinie\kreaturattribute{GE 14, KK 8, KO 8, MU 6}
    }}
\newcommand{\kreaturwolfrudelfuehrer}{\kreatur{Wolf Rudelführer}{verbreitetes Rudeltier}{gfx/kreaturen/tier}{
        \kreaturkampfwerte{6}{2}{8}{2}\trennlinie\kreaturvorteile{Angepasst II (Dunkelheit)}\trennlinie\kreaturwaffe{Biss}{0}{10}{10}{2W6}{}\kreaturkampfvorteile{Niederwerfen, Standfest, Sturmangriff, Haltet stand!, Keine Gefangenen!}\trennlinie\kreaturattribute{GE 14, KK 6, KO 6, MU 6}
    }}


\usepackage{listings}
\lstset
{
    language=[LaTeX]TeX,
    breaklines=true,
    basicstyle=\texttt\scriptsize,
    keywordstyle=\color{dunkelrot},
    identifierstyle=\color{braun},
}

\begin{document}

\hauptteil
\titelseite{%
    \begin{centering}
        \titelbild{bilder/template}
        \platz
        \titel{Noch ein Ilaris}\\
        \titel{Latex Template}
        \platz
        \platz
        \fanprodukt
    \end{centering}
}
%\kapitel{Credits}
\neueseite
\phantom{} % unsichtbares element, damit \platz am seitenanfang funktioniert
\platz
\credit{\LaTeX -Klasse}{\link{https://github.com/Ilaris-Tools/IlarisTex}{Lukas Ruhe}}
\credit{Illustrationen}{Bernhard Eisner}
\credit{Titelbild}{\link{https://unsplash.com/photos/Hi0bdO0vEfo}{Ivana Cajina}}
\credit{Glossar}{\link{https://github.com/XaverStiensmeier/IlarisGlossaryEntries}{Xaver Stiensmeier}}
\credit{Inspiriert durch}{\link{https://github.com/Ilaris-dev/Ilaris-LaTeX}{Ilaris-Latex}\\von Jan Battenfeld}
\lizenz
\inhaltsverzeichnis
%\hauptteil
\kapitelihvz{Das Vorwort}{Vorwort}
\spaltenanfang
\abschnitt{Latex Dokument}
Es gibt eine Grundstruktur, die jedes Dokument, das auf diesem Template aufbaut, befolgen sollte.
\begin{lstlisting}
    \documentclass{Ilaris}
    \begin{document}
    \anfang
    ...
    \hauptteil
    ...
    \end{document}
\end{lstlisting}
\neuespalte


\textquotedblleft Zwischen den Befehlen \gl{anfang} und \verb!\hauptteil! steht alles, was vor dem Inhaltsverzeichnis kommt und noch keine richtigen Seitenzahlen hat. Zum Beispiel Titelseiten, Impressum oder ein Vorwort. Optional kann nach dem Hauptteil noch ein \verb!\schluss! eingeleitet werden, in dem zum Beispiel Anhänge oder das Glossar untergebracht werden können. Für kurze schlichte Dokumente ist das jedoch oft nicht nötig.

\abschnitt{Digital vs Druck}
Aktuell ist das Template auf ein digitales Dokument ausgelegt. Es wurden weder leere extra Seiten für die Rückseiten des Covers eingefügt, noch Druckränder oder Tintensparende Hintergründe verwendet. Stattdessen, wurde Wert darauf gelegt, verschiedene Stellen (Inhalt, Glossar, Anhang, Handouts...) mit einander zu verlinken, sodass beim Lesen des Dokuments schnell an die richtigen Stellen gesprungen werden kann. Links sind unterstrichen hervorgehoben. Es kann natürlich ohne Weiteres ausgedruckt werden, wie es ist. Für zukünftige Versionen ist jedoch noch eine Option für eine extra Druckversion geplant.
\spaltenende

\kapitel{Layout}
\abschnitt{Seiten und Spalten}
Ein neues Kapitel wird mit dem Befehl \gl{kapitel} eingeleitet. Es wird automatisch formatiert und im Inhaltsverzeichnis registriert. Lange Kapitelüberschriften werden zwar automatisch umgebrochen, aber es sieht trotzdem besser aus die Überschriften kurz zu halten. Überschriften beginnen automatisch eine neue Seite, sollte ansonsten manuelle Seiten oder Zeilenumbrüche gebraucht werden, gibt es \textbf{\gl{neuespalte}} und \gl{neueseite}.
Ein Großteil, des Dokuments wird wahrscheinlich im Zweispaltenlayout geschrieben. Dieses wird mit \gl{spaltenanfang} begonnen und muss unbedingt mit \textbackslash Spaltenende beendet werden. Da die Kapitelüberschriften beide Spalten überspannen sollen, kann im Zweispaltenlayout \gl{skapitel} benutzt werden, um nicht jedes Mal das Spaltenlayout wechseln zu müssen.
\dreispaltenanfang
Ein dreispaltiges Layout kann entsprechend mit \gl{dreispaltenanfang} begonnen werden.\\
Ein dreispaltiges Layout wird im Regelbuch beispielsweise für Zauber und Liturgien verwendet.
\spaltenende

\spaltenanfang

\abschnitt{Titelseite}
Die Titelseite baut auf einem Bild auf, das idealerweise schon im DinA4-Format (Portrait) vorliegt. Erstellt wird sie mit dem Befehl \gl{titelseite}, wobei alles was auf die Titelseite kommt mit in die Klammern gehört. Ein paar vordefinierte Befehle, die nützlich sein könnten um die Titelseite einfach zu gestalten: \gl{titel}

\abschnitt{Überschriften}
Jedes Kapitel kann in einzelne Abschnitte (Unterkapitel) gegliedert werden, die mit \gl{abschnitt} eingeleitet werden und ebenfalls im Inhaltsverzeichnis auftauchen. Die Abschnittsüberschriften sollen das Spaltenlayout (falls verwendet) nicht unterbrechen.
\geschichte{Ich bin eine Geschichte oder ein Auszug aus einer Schrift, wie sie in der Spielwelt existieren könnte. Auch Zitate, Informationen aus Büchern oder Liedtexte könnten so aussehen.}{Autorin}

Diese kursiven Textblöcke, können mit dem Befehl \gl{geschichte} erstellt werden. Kurze Zitate oder einzelne \anf{Wörter} können mit dem Kurzbefehl \gl{anf} auch in Anführungszeichen gesetzt werden. Dieser Text hier ist ein ganz normaler Standardtext, so wie man ihn herunterschreibt ohne irgendwelche Layoutbefehle zu verwenden. Man kann ihn durch Überschriften in Abschnitte unterteilen und diese wiederumtel in Absätze unterteilen.

\absatz{Zauber- und Liturgieeigenschaften}
\textsc{Viele} Zauber und Liturgien verfügen über eine dieser
Eigenschaften, weswegen wir sie hier gesammelt haben.

\absatz{Der Absatz}
Beschriftete Absätze sind Unterunterüberschriften und stehen nicht mehr im Inhaltsverzeichnis, heben sich aber noch deutlich ab. Sie werden mit \gl{absatz} erstellt. Wenn nur der Textfluss abgesetzt werden soll, ohne irgendwelche Überschriften, geht das durch eine einfache Leerzeile im Text. Einfache Zeilenumbrüche werden von Latex ignoriert, können aber mit \gl{neuezeile} erzwungen werden. Die letzte Ebene der Struktur bildet der \gl{unterabsatz}.

\unterabsatz{Unterabsatz}
Hier steht dann was zu einem ganz bestimmten Thema oder für einen ganz bestimmten Ort. Im Inhaltsverzeichnis taucht das aber nicht auf...

\abschnitt{Kästen}
Kapitel fangen automatisch auf neuen Seiten an.
Ein paar auf die schnelle hingetüdelte Beispiele um Kästen und andere Umgebungen umzusetzen. Für einige dieser Kästen, hab ich das tcolorbox Paket ins Spiel gebracht, das nimmt einem etwas Arbeit ab, wenn man noch weitere oder kompliziertere Kästen erstellen will. Oder soetwas wie Seitenumbrüche oder transparente Hintergründe einstellen will (Kreaturenkasten). Für die Transparenz im Teil des Kastens wird allerdings eine relativ aktuelle Version des Pakets benötigt, falls die Installation zu Umständlich ist, kann man natürlich auch einfach ein untransparente passende Farbe nehmen.

\absatz{Pergamente}
Ein Pergamentkasten eignet sich wunderbar um einzelne Informationen deutlich hervorzuheben. Einen farbigen Kasten erstellt man mit \gl{kasten} und \gl{kastengrau} erstellt einen ähnlichen aber etwas graueren Kasten. Die können zum Beispiel Erzähltexte oder Informationen für gelungene Talentproben sein. Pergamente sind auch Rechteckig, auch wenn das folgende am Rand schon etwas ausgefranzt ist.

\kasten{
    Dies ist der Standardkasten. Da noch Bilder für Schmale, Breite und Größere vorliegen, mögen eventuell noch etwas speziellere Kästen folgen.\\
    Für kleine Absätze im Zweispaltenlayout sollte es aber erstmal reichen.
}

Kästen über die volle Breite können einfach außerhalb des Spaltenlayouts erstellt werden. Im Kapitel Handouts gibt es auch DinA4 Kästen an. Es können auch grauere Kästen verwendet werden, die zwar mehr aus einem \anf{Fehler} entstanden sind, aber vielleicht sind sie ja mal von Nutzen.

\kastengrau{
    Hintergrund dieses Kastens ist das Originalbild, das nicht mit dem beigen Hintergrund der Seite multipliziert wurde. Der Unterschied wird im Abschnitt Illustrationen veranschaulicht.
}

\spaltenumbruch

\absatz{Kreaturen}
Den Kreaturenkasten habe ich absichtlich erstmal nach einer allgemeinen Kreatur benannt und nicht einfach Gegner. Dafür gibt es zwei Gründe. Zum einen, kann es ja auch freundlich gesinnte Wesen geben und der selbe Kasten soll auch diese darstellen können, Kampfwerte oder freier Text für kommentare sollten dabei frei kombinierbar und optional bleiben. Zum anderen, kann (sofern erlaubt), der Befehl für Gegnerkästen noch ergänzt werden um mit existierenden Ausgaben von Konvertierungstools etc, kompatibel zu bleiben.

\kreatur{Name}{Beschreibung}{gfx/kreaturen/tier}{
    \kreaturinfo{Info}{Text}
    \kreaturinfo{Kampfvorteile}{Standfest, Sturmangriff}
    \kreaturinfo{Fertigkeiten}{Eins zwei drei und noch mehr text damit man die}
    \linie
    \kreaturinfo{Fertigkeiten}{Eins zwei drei und noch mehr\\ text damit man die}
}

\kreatur{Mensch}{Humanoid, Unfreundlich}{bilder/zwerg}{
    \kreaturinfo{Vorteil}{Er muss sich nicht ständig unterhalten, weil ihn eh keiner mag}
    \trennlinie
    \kreaturinfo{Nachteil}{Ihm ist oft langweilig, weil sich niemand mit ihm unterhalten möchte}
}

Ich finde die Eigenschaften blockweise einzuschieben wesentlich einfacher und flexibler, als einen Befehl mit fest vorgegebenen Argumenten zu benutzen. So können die Kreaturen oder NSCs so auch um beliebige Felder, wie zB möglichen Loot oder Zitate eines NSCs ergänzt werden.



\absatz{Proben}
%\scalebox{-1}[1]{\includegraphics[width=\linewidth]{layout/grafik/proben/verschlagenheit}}
%\includegraphics[width=\linewidth]{layout/grafik/proben/alchemie}
Für Probenboxen kam mir die Idee, die die Fertigkeiten hübschen Fertigkeiten-Bildchen zu verwenden. Die aktuelle Umsetzung mit dem gebastelten Rahmen ist vielleicht etwas zu fancy und hat den Nachteil, dass ie nicht wirklich in der Höhe skalierbar ist ohne das Bild zu verzerren. Alternativ gibt es ja auch schon die Möglichkeit die Pergamentkästen für Proben zu verwenden.

\probe{wahrnehmung}{Infanteriewaffen und\\ Speere (III, 18)}{Ich bin eine Probe. Die Box ist leider nicht skalierbar, daher müssen die Erzähltexte kompakt oder stichpunktartig sein, oder ein anderer Kasten gewählt werden.}

\probe{magiekunde}{Magietheorie (III, 21)}{%
    \tabelle{rX}{
        +8 & Doppelte Schwierigkeit \\
        +4 & Schwieriger    \\
        0 & Normal \\
        -4 & Vereinfacht \\
        -8 & Doppelte Vereinfachung \\
    }

}

\kasten{
    \kastenkopf{Magietheorie (III, 21)}
    \tabelle{rX}{
        +8 & Doppelte Schwierigkeit, dessen Ergebnistext etwas mehr als ein oder zwei Zeilen einnimmt \\
        +4 & Schwieriger, hat auch informationen mit etwas mehr Text    \\
        0 & Normal, nur kurz \\
        -4 & Vereinfacht, wieder etwas mehr Text um zu sehen wie der Kasten sich strecken lässt \\
        -8 & Doppelte Vereinfachung \\
    }
}

\absatz{Info}
Für kleine Anmerkungen oder als Alternative zu den Pergamentkästen gibt es noch kleine Infokästen.
\info{
    Hier kann alles mögliche drinstehen, zum Beispiel Spoilerwarnungen, Verweise auf andere Abenteuer oder sinnloser Text wie dieser.
}

Und, wenn auch nicht wirklich als Kasten erkennbar, kann man einzelne Textpassagen mit einem beliebigen Symbol versehen. Dies könnte zum Beispiel für Beschreibungen von NSCs zusammen mit ihrem Porträt benutzt werden, oder mit kleinen wieder kehrenden Symbolen, die dem Leser direkt eine Zugehörigkeit zu einem bestimmten Thema signalisieren: Zum Beispiel, durch das Wappen einer Fraktion oder einem Symbol für Geheimnisse, Fallen usw.

\absatz{Symboltext}

\symboltext[2cm]{bilder/einleitung.jpg}{
    Das Symbol sitzt links am Beginn des Absatzes. Der Text umfließt das Symbol automatisch. Die Breite des Symbols kann als optionales Argument übergeben werden. Per default steht es auf 1cm. Es kann sich lohnen etwas mit der Größe zu variieren, wenn die Zeile darunter zu viel/wenig Abstand hat.
}

\absatz{Stichpunktlisten}
Aufzählungen haben jetzt auch ein eigenes Symbol:

\stichpunktliste{
    \punkt Der Befehl in der Klasse wäre \em{stichpunktliste}
    \punkt Die klassische itemize-Umgebung sollte aber genauso funktionieren.
}

\absatz{Karten}
Wer es mag sich kleine Karten für Manöver, Zauber, Regeln oder NSCs auszudrucken kann diese Kästen benutzen. Sie sind haben die Maße von üblichen Spiel- und Sammelkarten (63mm x 88mm, zB Magic). Und sollten damit in handelsübliche Sleeves und Sammelmappen passen. Dabei muss natürlich beachtet werden, dass die Karten ggf. minimal kleiner werden, wenn die Seite beim Drucken herunterskaliert wird.

\karte{test}


\abschnitt{Tabellen}
Tabellen können leider noch nicht vollständig durch Kurzbefehle ersetzt werden, dafür stehen zu viele Variablen im Raum. Die Anzahl, Ausrichtung und Größe von Spalten und Zeilen sollen variabel bleiben, eine Kopfzeile optional. Das Streifenlayout wird allerdings automatisch angewendet wenn eine \gl{tabelle} begonnen wird. Der erste Eintrag enthält dabei das Spaltenlayout, das im Kapitel Latex Tipps auf Seite \pageref{latextipstabellen}
genauer erklärt wird. Eine ganze Zeile zu formatieren, ist etwas komplizierter, daher als Übergangslösung erstmal die Möglichkeit jede Zelle der Kopfzeile mit \gl{tkopf} zu formatieren.

Das zweite Argument gibt den Inhalt der Tabelle direkt an Latex weiter.Da muss man zur Zeit noch echten Latexcode schreiben, weil mir auf die schnelle kein Praktischer Kurzbefehl eingefallen ist. Aber da kann man sich sicher noch etwas überlegen, um das deutlich zu verinefachen. Eigentlich steht der Text hier auch nur um die Tabelle im Abschnitt nicht so alleine zu lassen.

\tabelle{X X X}{
    \tkopf{eins} & \tkopf{Zwei} & \tkopf{Drei} \\\hline%
    ich      & bin      & eine     \\
    einfache & latex    & tabelle  \\
    in       & tabularx & umgebung \\
}


\neueseite
\abschnitt{Illustrationen}


Illustrationen auf weißem Hintergrund, die sich dem Seitenhingergrund anpassen sollen, können (multiplikativ) über die Seite blenden. Dabei entsteht der Effeckt, als würde direkt auf das pergament gezeichnet/gedruckt worden sein. Abbildungen dieser Art werden hier Zeichnungen genannt.
\zeichnung{bilder/einleitung}
Sollen die Illustrationen in ihren echten Farben dargestellt werden, nennen wir es Bild und es würde so aussehen
\bild{bilder/einleitung.jpg}

Um Bilder nicht in voller Größe anzuzeigen, kann als optionales Argument die gewünschte Breite in mm oder cm angegeben werden:

\czeichnung[3cm]{bilder/einleitung}

\zeichnung[3cm]{bilder/einleitung}
\bild[3cm]{bilder/einleitung}


\begin{centering}
    \zeichnung[3cm]{bilder/einleitung}
\end{centering}

Um echte Farben zu erhalten aber trotzdem keinen weißen Rand zu haben, kann man natürlich auch Bilder mit transparentem Hintergrund (.png) verwenden.
Bilder können auch in einen Zierrahmen eingefasst werden. Das sieht zum Beispiel so aus:
\rahmen{bilder/einleitung}

Um Bilder oder Zeichnungen über die volle breite darzustellen, kann das Spaltenlayout einfach kurz unterbrochen werden (spaltenende, spaltenanfang).
\spaltenende
\rahmen{bilder/einleitung}
\zeichnung{bilder/einleitung}
\kapitel{Großes Bild mit Überschrift}
% \makebox[\linewidth][c]{%
%     \begin{minipage}{15cm}
%         \rahmen{bilder/template}
%     \end{minipage}}%
\rahmenklein[15cm]{bilder/template}
\spaltenanfang


\skapitel{Handouts}
Handouts sind für viele Spielhilfen und Abenteuer keine Seltenheit und niemand hat Lust irgendwelche ausgeschnittenen oder halb abgedeckten Seiten des ausgedruckten Dokuments auf den Spieltisch zu legen. Deshalb gehören Handouts in den Anhang. Um möglichst keine unnötige Informationen mit auf den ausgedruckten anhängen zu haben, wurde ein extra Seitenlayout erstellt, mit anderen Seitenbegrenzungen, ohne Hintergrundbild Seitenzahlen etc. Die zwei arten von Handouts die mir am häufigsten untergekommen sind, sind wohl Text (Notizen, Briefe, Buchseiten etc..) und Karten. Für letzteres (und ersteres mit besonderen Designwünschen) bietet sich ein Bild für die ganze Seite an. Für ersteres eventuell die schon verwendeten Kästen in der richtigen Größe.
Ein erster Ansatz für ein Handout kann schon auf Seite \pageref{beispielkasten} betrachtet werden. In der nächsten Version könnte aber noch einiges an Komfort und weitere Beispiele folgen.


\skapitel{Extras}
Dieses Kapitel beinhaltet Funktionen, die nicht zur Kernklasse gehören, aber durch \textbackslash include{extras} und ggf. weiterer \anf{Plugins} aktiviert werden können. Die Idee dieser optionalen Module ist, sie so einfach wie möglich nutzbar zu machen, ohne dass sie dabei die eigentliche Klasse mit ggf. unnötigen Paketen und Einstellungen verkomplizieren. Auch soll es so für andere Entwickler möglich sein, die Klasse einfach erweitern zu können. Aktuell gibt es zB die Möglichkeit auf eine Sammlung von Kreaturen und von Glossareinträgen zurück zu greifen.

\abschnitt{Glossar}
Die IlarisGlossaryEntries sind mit eingebunden und können direkt benutzt werden. Damit können Erklärungen zu Regeln, Zaubern, Manövern und vieles mehr automatisch im Anhang generiert und für die schnelle Nutzung am PC an entsprechenden Stellen verlinkt werden. Die voreingetragenen Texte werden über die Zeit erweitert. Ein Beispiel wäre der Vorteil:
\gl{angepasstI_Vorteil}.
Natürlich können auch eigene Glossareinträge für das jeweilige Dokument erstellt werden:
\glossareintrag{glossareinträge}{name={Glossar}, description={Das Glossar ist eine Liste von Begriffen, die am Ende des Dokuments erklärt werden. Der entsprechende Begriffe kann im Dokument angeklickt werden um direkt zu dieser Erklärung zu springen. Die Zahlen hinter der Erklärung geben an wo (Seitenzahl) die Begriffe vorkommen und können ebenfalls angeklickt werden um zurück zu springen.}}
\gl{glossareinträge} eingetragen werden.
Für dieses Template wird das Glossar vorallem dazu verwendet die Latex-Befehle zu erklären.

\spaltenumbruch

\abschnitt{Kreaturendatenbank}
Testweise habe ich mal eine Datenbank für Kreaturen als yaml-Datei angelegt, die automatisch in eine tex-Datei exportiert werden kann. Das Ziel ist, diese DB gemeinsam zu erweitern und in Latex können die Blöcke für jede Kreatur in der Sammlung automatisch erzeugt werden. Die Autovervollständigung, kann so beim Schreiben sogar noch als Inspiration für spontane Gegner dienen. So ergbit der Befehl \textbackslash kreaturelefant zum Beispiel folgenden Kasten, ohne das irgendwelche Werte angegeben werden müssen:

\kreaturelefant
\kreaturbaumdrache
\skapitel{Latex Tips}

\abschnitt{Tabellen}
\label{latextipstabellen}
TODO: Spaltenlayout bei Tabellen
\spaltenende

\glossar
%\schluss
\handout




\kastenlink{beispielkasten}

\kastenlink{kasten2}

\handoutkasten{beispielkasten}{
    {\huge\handschrift
            \kastenname[Beispielkasten - Linktest]{beispielkasten}
            Das hier ist ein Pergamentkasten, der benutzt wird um direkt im Latex Dokument ein Handout zu erstellen.
            Man könnte hier natürlich noch eine Schriftart wählen, die das ganze nach einer krakeligen Handschrift oder feinster Kaligraphie aussehen lässt, je nachdem, wofür das Handout benutzt werden soll.\\[1cm]
            Der \gl{handout} Befehl beginnt eine neue weiße Seite, auf der sich ausgetobt werden kann. Alternativ gibt es diesen vorgefertigten Kasten (\gl{handoutkasten}) und Handouts für Bilder (\gl{handoutbild}) oder um direkt fertige PDFs einzubinden (\gl{handoutpdf}).}
}
\handoutkasten[8cm]{handoutnotiz}{
    \kastenname{kasten2}
    {\Large\krakelschrift
        Oder vielleicht ein kleiner Zettel mit einer krakeligen Notiz, der gefunden werden kann.
    }
}
\neueseite%
%\begin{multicols}{3}
\karte{\textbf{Angepasst (Schnee)\\}\glsdesc{angepasst(Schnee)I_Vorteil}}%
\karte{Zweite Karte}%
\karte{Karte Drei}\\%
\karte{Vierte Karte}%
\karte{test}
%\end{multicols}


\handoutende
\end{document}